\section{Introduction}\label{sec:intro}
% =====================================================================
% Section 1: Introduction.
% Sources: research/00-brief.md (questions, verbatim); the final records research/tracks/*/notes-final.md
% (results in brief, caveats, open problems) and their referee reports; research/prior/.
% Every claim below points to the section result that states it with its hypotheses.
% Final pass after the paper review (research/paper-review/edits/E1-results.txt, E2-results.txt, DECISIONS.md):
% omega-step wording (holds in every model of true arithmetic), "not derivable in general", (Rich) and (Rich_c)
% qualifiers on "only if" anchor claims, ReplS truth-soundness only for the lgg of anchor data, capture-guard
% corrections, separation proved for PA and (ideal refuter) the experimental mixtures but only sampled for ZF/ZFC,
% sharing pass guaranteed only for targets labelling a phase-one cluster; crisp answers added in sec:intro:answers.
% =====================================================================

A mathematical community uses axioms: single sentences, such as Extensionality, and schemas, such as induction in Peano
arithmetic or Separation in set theory, each standing for infinitely many sentences. Its proofs show only instances,
often without saying which axiom they instantiate. We ask when a learner can recover from such positive instances
exactly which axioms the community uses. The aim is checking, not proving: we want a verifier that accepts a candidate
axiom step iff it is an instance of an axiom in use, and that stays sound when the candidates come from an untrusted,
possibly adversarial, prover.

This is the setting of the report \citet{claude2026whatfollows} (\emph{the parent report}) on learning to check
inferences. Its central tool is the \emph{cautious} verifier: fix a class of hypotheses and accept a sentence iff every
hypothesis in the class that contains the data contains it. If the target is a member of the class (more precisely, an
intersection of members), this verifier is sound at all times against every prover, and it is exact once the data
contain an \emph{anchor}, a finite set of instances that every covering hypothesis extends to the whole target
(\cref{sec:setting:learning}). For rules cited by name the parent report used first-order anti-unification
\citep{plotkin1970note,reynolds1970transformational}. PA induction does not fit: $\Ind(\varphi)=(\varphi(0)\wedge\forall
x(\varphi\to\varphi[Sx/x]))\to\forall x\varphi$ substitutes terms into the motive $\varphi$. A follow-up analysis
(\emph{the earlier analysis}; unpublished, its record is in the accompanying repository, and we re-prove what we use)
learned induction with \emph{determinate} second-order templates. In these every metavariable has an occurrence applied
to distinct bound variables, such as $P(x)$ in $T_{\Ind}=P(0)\wedge\forall x(P(x)\to P(Sx))\to\forall x\,P(x)$, and
arguments contain no metavariables (\cref{rem:setting:nested} shows why). Matching is unique and linear, and two
instances are an anchor iff their motives have different main symbols and at least one has $x$ free. Kaarel H\"anni then
asked three questions about this method.

\subsection{The questions}\label{sec:intro:questions}

Verbatim:
\begin{quote}
\begin{enumerate}[label=\arabic*.,leftmargin=1.5em]
\item does the same thing work for learning an axiom of the form forall x phi(x) from a bunch of sentences of the form
  phi(x)?
\item does the same thing work for learning each axiom schema of ZFC from a bunch of instances?
\item anyway can you try to come up with a method that works for all these? ideally, your method would also work for
  learning many axiom( scheme)s at once, even if they do not come with labels for instances of which axiom scheme they
  are
\end{enumerate}
\end{quote}
``The same thing'' is the method for induction. We call these Questions 1--3 (not to be confused with the axioms Q1--Q7
of Robinson arithmetic $\Q$). We read the data of Question~1 as instances $\varphi(t)$ at terms $t$: since all data are
read under universal closure, a datum $\varphi(x)$ with $x$ free would already be the axiom (\cref{prop:univ:gen}).

The answers compare four classes of templates (\cref{sec:setting:templates}): first-order patterns $\FO$, whose
metavariables are $0$-ary; Miller's higher-order pattern templates $\PAT$ \citep{miller1991logic}, in which every
metavariable occurrence is applied to distinct bound variables; determinate templates $\DT$, in which every metavariable
has at least one such occurrence and the others may have any metavariable-free arguments, as $P(0)$ and $P(Sx)$; and
rigid second-order templates $\SO$ without this requirement. So $\FO\subseteq\PAT\subseteq\DT\subseteq\SO$. Data are
formulas whose free variables are parameters, read under universal closure; bound variables are de Bruijn indices, and
values of metavariables contain no free bound variable, so capture cannot occur (the $\lambda$ convention). Besides
data, a learner may use a \emph{refutation oracle}: sound but one-sided evidence that a sentence is false, such as a
$\Delta_0$ counterexample in $\N$ or a hereditarily finite counterexample to a $\Pi_1$ sentence of set theory
(\cref{sec:setting:refutation}).

\subsection{The answers in brief}\label{sec:intro:answers}

\begin{itemize}[leftmargin=1.5em]
\item \emph{Question~1: yes for the instance schema, not for the sentence.} The instance schema $\varphi(z)$ is a
  first-order pattern; two instances whose terms have different head symbols are an anchor (with several variables,
  Plotkin's distinctness is needed too), for first-order anti-unification and the determinate learner alike. Passing
  from the closed instances to $\forall x\varphi$ is an $\omega$-rule step: truth-safe for every $\varphi$ exactly when
  the closed-term substructure is elementary (as in $\N$, not in $\R$), and not derivable in general. A cautious learner
  takes this step silently when parameters are admissible values; a learned closedness guard prevents it, and data with
  a parameter at every variable license it as ordinary generalization.
\item \emph{Question~2: yes, and more easily than for induction.} The axioms other than Separation and Replacement are
  single sentences, each its own anchor. Every ZF schema, in each formulation considered, is a Miller pattern, and in
  every class from $\PAT$ to $\SO$ finitely many instances are an anchor iff their bodies' main symbols are not all
  equal and every argument place is used by some body; two instances can suffice. Pattern anti-unification and the
  determinate learner recover each schema from any anchor. Plain first-order anti-unification works only for some
  formulations and encodings. PA induction is not a pattern; first-order and pattern anti-unification fail on it, and
  the determinate learner learns it.
\item \emph{Question~3: yes, under a condition on the world that we make explicit.} The method \DTRC{} clusters untagged
  data by refuting the templates that would merge them, then verifies each cluster cautiously. If every template
  covering data of two targets has a refuted instance (\emph{refutation separation}), it recovers the hidden labels, is
  target-sound at all times and is exact once each target's data contain an anchor, so it needs no more data than a
  learner given the labels. Separation is proved for $\Q$'s axioms with induction and checked on sampled ZF and ZFC
  data. Without it \DTRC{} is sound only relative to a residue of unrefuted merges, and no computable learner can avoid
  depending on the refutation depth.
\end{itemize}
The rest of this subsection gives the results behind these answers; \cref{tab:intro:anchors} collects the anchor
conditions.

\begin{table}[t]
\centering\small
\setlength{\tabcolsep}{4pt}
\begin{tabularx}{\textwidth}{@{}>{\raggedright\arraybackslash}p{0.2\textwidth}
  >{\raggedright\arraybackslash}X >{\raggedright\arraybackslash}p{0.15\textwidth}
  >{\raggedright\arraybackslash}p{0.15\textwidth} >{\raggedright\arraybackslash}p{0.17\textwidth}@{}}
\toprule
target & $\FO$ (first-order lgg) & $\PAT$ (pattern lgg) & $\DT$ & $\SO$\\
\midrule
single axiom ($\Ext$, $\AC$, \dots) & \multicolumn{4}{l}{the axiom itself, in every class}\\
instance schema of $\forall\bar x\varphi$ & \evR+\evD & \evR+\evD & \evR+\evD & \evR+\evD\ do not suffice\\
ZF schema & only some formulations; most others: no sound finite union & \evR+\evN & \evR+\evN & \evR+\evN\\
PA induction & unsound; no sound finite union & not in $\PAT$; pattern lgg unsound & \evR+\evN & \evR+(B)\\
any $T^*\in\DT$ & & & \evRs+\evN+\evU & \\
\bottomrule
\end{tabularx}
\caption{When finitely many genuine instances are an anchor, by target and class (with the class's anti-unification
learner where there is one). \evR: the root symbols of the values (terms, bodies, motives) are not all equal; \evD:
Plotkin's distinctness; \evN: every argument place is used by some value; (B): some motive has a free occurrence of $x$
inside no compound term whose only variable is $x$; \evRs, \evU: root variation at every occurrence and no coincidence.
``If'' directions need no assumption; some ``only if'' directions use mild richness assumptions (\Rich, or enough closed
terms). Sources, by row: \cref{prop:zf:ground}; \cref{thm:univ:anchor,prop:univ:dt} (the $\PAT$ entry follows from
$\FO\subseteq\PAT\subseteq\DT$; the exact $\SO$ condition is open); \cref{tab:zf:classify,thm:zf:nounion,thm:zf:anchor};
\cref{prop:zf:indtwo,thm:zf:indunion,prop:zf:indpat,thm:zf:indanchor,prop:zf:indso}; \cref{thm:single:anchor}.}
\label{tab:intro:anchors}
\end{table}

\paragraph{Question 1 (\cref{sec:univ}).}
Plotkin's events are \evR: the terms at each $x_i$ do not all have the same root symbol, and \evD: for $i\neq i'$ some
datum puts different terms at $x_i$ and $x_{i'}$. Closed instances satisfying them are an anchor for the instance schema
$\sigma_\varphi=\varphi(z_1,\dots,z_k)$ in $\FO$ and in $\DT$, and if the language has enough closed terms, as $L_A$
does, only these are (\cref{thm:univ:anchor,prop:univ:dt}); in $\SO$ the events do not suffice. For one variable the
probability that $N$ i.i.d.\ instances violate \evR\ is exactly $\sum_fp_f^N$, $p_f$ the probability of root $f$
(\cref{thm:univ:rates}). The $\omega$-step is truth-safe in $M$ for every $\varphi$ iff the closed-term substructure of
$M$ is elementary (\cref{prop:univ:tv}, by the Tarski--Vaught criterion \citep{tarski1957arithmetical}). This holds in
$\N$ and in every model of true arithmetic, but not in the ordered field $\R$, where $\neg(x\cdot x=1+1)$ is true at
every closed term, nor in the universe $V$ of sets or in a model of $\PA+\neg\Con(\PA)$ (\cref{ex:univ:fail}). In
general the step is not derivable: $\Q\vdash0+\num n=\num n$ for each $n$, but $\Q$ with all these instances does not
prove $\forall x\,(0+x=x)$ (\cref{ex:univ:q}). With parameters admissible, the cautious learner in $\FO$ or $\DT$
accepts $\varphi(\bar p)$, which is $\forall\bar x\varphi$, after the first anchor (\cref{prop:univ:open}). For
first-order schemas a closedness guard learned from closed data prevents this ($\DT$ has no guards); data with a
parameter at every variable make the step ordinary generalization (\cref{prop:univ:gen}). If a metavariable may take a
bound variable (named variables, or de Bruijn indices treated as constants), the schema can gain capture instances, such
as $\exists y\,\neg(y=y)$ from $\exists y\,\neg(y=x)$, false in every structure; whenever capture instances exist, the
unguarded verifier accepts some after an anchor. Learned guards repair this: closedness excludes capture as well, and a
no-capture guard keeps it excluded once parameter data have removed closedness (\cref{prop:univ:capture},
\cref{tab:univ:learner}).

\paragraph{Question 2 (\cref{sec:zf}).}
The ZF schemas, in each textbook formulation considered (\cref{tab:zf:forms}), are patterns because $L_\in$ has no
function symbols: an instance only renames variables inside the body (\cref{thm:zf:classify}). Their anchors are
characterized in \cref{thm:zf:anchor}; pattern anti-unification \citep{pfenning1991unification,baumgartner2017higher}
and the $\DT$ learner recover each schema from any anchor (\cref{cor:zf:patlgg}), with exact failure probabilities
(\cref{thm:zf:rates}). First-order anti-unification works iff the metavariable is applied to literally the same
variables at every occurrence, which depends on the formulation and on how bound variables are written
(\cref{thm:zf:classify,prop:zf:alpha}; \cref{tab:zf:classify} in \cref{app:zf}). Separation passes, with a freshness
guard learned from the data; without it the lgg has Russell's instance (\cref{prop:zf:russell}). In the other failing
cases with textbook names or de Bruijn indices, no finite union of guarded first-order schemas covers the schema without
a false member (\cref{thm:zf:nounion,prop:zf:false}). The one exception is spelled-out Replacement with textbook names:
the lgg of anchor data, with the guards learned from them, over-generalizes truth-soundly, as all its guarded instances
are theorems of ZF, and of a base theory exactly when that theory proves Collection; no refutation can expose it
(\cref{prop:zf:coll}).

PA induction is not a pattern. On raw instances first-order anti-unification can be unsound from two instances and is
unsound on every anchor, no finite union of first-order schemas is truth-sound, some unsound lggs are never refuted by
true closed quantifier-free sentences, and pattern anti-unification is unsound too
(\cref{prop:zf:indtwo,thm:zf:indunion,thm:zf:K0,prop:zf:indpat}). The $\DT$ learner is exact under \evR+\evN\
(\cref{thm:zf:indanchor}), and $T_{\Ind}$ is the least of the 26 $\DT$ templates covering any anchor
(\cref{prop:zf:indgen}). The obstruction is the substitution $\varphi[Sx/x]$, not induction: recording it, or Tarski's
substitution-free form $\IndEq$ \citep{tarski1965simplified}, makes induction a first-order pattern, and $\Q$ plus the
instances of $\IndEq$ axiomatizes $\PA$ (\cref{prop:zf:sub,prop:zf:indeq}).

\paragraph{The general theory (\cref{sec:single}).}
These anchor conditions are cases of one theorem. For a target $T^*\in\DT$, finite data have finitely many minimal
covering templates, one of them below $T^*$ (\cref{thm:single:min}). The cautious verifier is a polynomial-time test of
``features'' of the data (\cref{thm:single:feat}), although there can be $4^n$ minimal templates
(\cref{prop:single:blowup}). Root variation at every occurrence, non-vacuity and no coincidence make data an anchor;
under a richness assumption \Rich, which always holds with parameters, every anchor satisfies them and its least
covering template is the target (\cref{thm:single:anchor,cor:single:anchor-lgg}). For first-order targets these events
are \evR+\evD, for one formula metavariable \evR+\evN\ (\cref{cor:single:special}). Failure probabilities decay
exponentially (\cref{thm:single:rates}), and the worst-case number of escalations (queries passed to an oracle) is
$\Theta(N^2)$ in the sentence size, which refutes an earlier conjecture that it is linear
(\cref{thm:single:esc,prop:single:quadratic}).

\paragraph{Question 3 (\cref{sec:many}).}
From positive data a bound on the number of targets is necessary: without one the cautious union learner accepts only
the data, whatever sound negatives it has (\cref{prop:many:bound}; \citealp{gold1967language}). With a bound, spare
slots cost data diversity that no sound negative can buy, because splitting a target into sound specializations is never
refuted (\cref{thm:many:splits,thm:many:slots}). With the exact bound and negatives that refute every template covering
anchor data of two targets, untagged learning is as easy as tagged learning, and the learner can find such negatives
itself (\cref{thm:many:pigeonhole}). \textup{\DTRC} (Determinate Templates with Refutation Clustering) needs no bound.
Starting from singletons, it repeatedly merges two clusters whose union has a minimal covering $\DT$ template with no
instance refuted within depth $d$, verifies each final cluster cautiously against its unrefuted minimal templates, and
asserts the union. Within-target merges are never refuted; so under separation the clusters are the hidden label classes
for every merge order, and the guarantees stated above follow, with the tagged learner's failure probability if
separation holds for all instances (\cref{thm:many:dtrc}). Refutation does the work that citations do in tagged data.

Without separation \DTRC{} is sound relative to a residue of unrefuted cross-target merges, truth-sound if these are,
and exactness can fail by fragmentation (\cref{thm:many:residue}); on some data no learner with a target-blind oracle is
both exact for one valid practice and target-sound for another (\cref{prop:many:ambiguity}). Learning $\forall x\varphi$
from its instances is such a merge, truth-sound iff $\forall x\varphi$ is true (\cref{prop:many:forall}). With
$\Delta_0$ evaluation in $\N$ and quantifier-free $\varphi$, ``refuted at no depth'' means ``true''
(\cref{prop:many:qf}); at a fixed depth this holds only relative to the residue, and for quantified $\varphi$ it can
fail. No computable learner is both sound relative to the full-depth residue and eventually exact on every practice
separated at some finite depth (\cref{thm:many:depth}, via the MRDP theorem \citep{matiyasevich1993hilbert}). The exact
$k$-union verifier is polynomial for $k\le2$ and coNP-complete for each fixed $k\ge3$; the refutation test of \DTRC{} is
NP-complete, and polynomial under \Rich\ once the cluster contains an anchor (\cref{thm:many:complexity}; not
independently refereed). On natural targets, $\Delta_0$ evaluation in $\N$ provably separates $\Q$'s axioms and
induction; hereditarily finite counterexamples and logic, with Russell's paradox for the comprehension-shaped axioms,
separate ZF and ZFC on every sampled cross pair and in every run (\cref{sec:many:examples}); and \cref{sec:exp} proves
separation of its PA and ZF mixtures for an ideal refuter (\cref{cor:exp:mixsep}). Weak forms of Union and Power merge
into a false universal-set schema that neither logic nor hereditarily finite counterexamples refute, only coherence
(\cref{prop:many:univset}). Minimum
description length is no substitute for refutation: it tracks the statistics of usage, not the logical boundaries of
schemas (\cref{sec:many:mdl}).

\paragraph{Experiments (\cref{sec:exp}).}
An implementation of the learner (in the class $\DTF$, with $0$-ary term metavariables) and of \DTRC{} matches the
predicted rates on Question~1 (mean $3.27$ instances to exactness, $3.25$ predicted). On Separation, a Replacement
variant, $\in$-induction and PA induction it is exact in most runs from $2$ or $3$ tagged instances, with no
non-instance accepted in $927$ probes (\cref{tab:exp:e2}). On mixtures of the PA and of the ZF axioms it recovers the
hidden labels in every run. With a pass that shares ambiguous data between clusters it matched a tagged learner given
every label a datum deserves, target by target, in every run (\cref{tab:exp:e3}); this is guaranteed only for targets
that label a cluster of the first phase, as all did here (\cref{prop:exp:share,ex:exp:share}). The limits show too: the
evidence for separation beyond the mixtures is narrow (\cref{tab:exp:e9}); similar targets merge into valid templates
(\cref{lem:exp:valid}); and with injected mistakes and refuter budgets of $80$ and $400$, \DTRC{} accepted a false
sentence, whose merge only a budget of $2000$ refuted (\cref{prop:exp:qF}).

\subsection{Contributions}\label{sec:intro:contrib}

Besides the answers, the paper gives a general theory of learning one determinate template (\cref{sec:single}); anchor
conditions for several metavariables, which differ between $\PAT$, $\DT$ and $\SO$ and are exact for the first two
(\cref{thm:zf:multi}); the method \DTRC{} with its theory of bounds, separation, residues, depth, noise and cost
(\cref{sec:many}); an implementation with baselines, a proof of separation for its mixtures with an ideal refuter, and a
reported failure (\cref{sec:exp}); and corrections to earlier claims of this project, such as the conjecture of linear
escalation bounds (\cref{app:ver}).

\subsection{What is not achieved}\label{sec:intro:limits}

\begin{itemize}
\item \emph{Separation is assumed, not guaranteed.} It is proved for $\Q$'s axioms with induction and, for an ideal
  refuter, for the experimental mixtures, and only computed on sampled ZF and ZFC pairs. It fails for some similar
  targets (\cref{lem:exp:valid,prop:many:univset}), no computable learner can know the separating depth
  (\cref{thm:many:depth}), and a characterization of separation in terms of the targets is open.
\item \emph{Refutation is one-sided and incomplete.} It never certifies that a template is the target
  (\cref{lem:setting:onesided}). Truth-sound over-generalizations are invisible to it (\cref{prop:zf:coll}), some
  unsound templates are never refuted by closed truths (\cref{thm:zf:K0}), and where the $\omega$-step is not
  truth-safe, a false universal can have only true closed instances, which no evaluation of closed instances refutes
  (\cref{prop:univ:tv,prop:univ:refute}). Our oracles are incomplete and the implemented refuter has a budget, so ``not
  refuted'' certifies nothing.
\item \emph{Cost.} In the worst case there are exponentially many minimal templates and quadratically many escalations,
  the union verifier is coNP-complete for $k\ge3$ and the refutation test NP-complete. The implementation uses $\DTF$
  without the feature verifier, on small synthetic data, not a corpus from a proof library.
\item \emph{Mistakes.} The theorems assume a valid practice and correct data. Unrefutable mistakes can be absorbed into
  a cluster or fragment it (\cref{prop:many:noise}), a frequent coherent family of mistakes cannot be told from a rule
  by positive data, and mistakes made the implementation accept a false sentence (\cref{prop:exp:qF}).
\item \emph{Rare targets and scope.} A target seen once is never generalized, a single axiom is learned only if it
  occurs among the data, and an unrefuted singleton mistake looks like a single axiom (\cref{sec:exp:limits}). We learn
  which instances a community uses, not the theory they generate, and from closed instances the truth of $\forall
  x\varphi$ is certified only as far as the oracle reaches. Open problems remain in every part; \cref{sec:disc} collects
  them.
\end{itemize}

\subsection{Reader's guide}\label{sec:intro:guide}

\Cref{sec:setting} fixes the framework and lists the named conditions (\cref{tab:setting:glossary}). \Cref{sec:single}
is the theory of one determinate template, used by the later sections. \Cref{sec:univ,sec:zf,sec:many} answer Questions
1--3, each ending with a summary that can be read first (\cref{sec:univ:summary,sec:zf:summary,sec:many:summary}).
\Cref{sec:exp} reports the experiments, and \cref{sec:disc} discusses the answers, related work and open problems. The
appendices, one per section, hold the full proofs and the secondary results the main text points to (numbered with a
letter, as \cref{prop:zf:indgen}); \cref{app:ver} says how the results were checked and what was corrected. Theorem-like
statements carry a status (\textsf{proved}, \textsf{computed}, \textsf{known}, \textsf{conjecture}, or a combination)
and a gray tag naming the result in the research record of the accompanying repository, which also holds the scripts and
the reports of adversarial referees; refuted conjectures stay labelled as refuted.
